\begin{beginnerstatement}[In simple terms: Reproducible results]

A reproducible build should produce the same result again from the same
inputs. For this case study, however, the evidence only shows that two
generations of \texttt{desktop-cloud + postgres} on \texttt{1cdfb6e} were
byte-for-byte identical in the same environment. A pinned dependency version
specifies exactly which version of a required library is used. A lockfile pins
such versions instead of selecting arbitrary newer packages on the next run. A
hash is a short digital fingerprint that can verify specific content very
reliably. The toolchain comprises the runtimes, compilers, and build tools in
use. Differences in the operating system, toolchain, network, or platform SDKs
can nevertheless affect reproduction.

For the Rust analyzer, provenance binds the exact compiler, target and
dependency versions and checksums for the build script, Cargo files, toolchain
file, and Rust sources to the checked-in WASI artifact. The build replaces
private source, home, Cargo, toolchain, and dependency paths with fixed virtual
paths. Three relocated Linux builds --- including paths with spaces and a local
vendor directory --- were byte-identical. This is narrow evidence for the
pinned Linux builder, not a guarantee of byte identity on macOS, Windows, or
with future compilers.

\end{beginnerstatement}
